\subsection{Comprobación de Nulos}

\subsubsection{Regla: referencias anulables y validación de argumentos}

\textbf{Regla:} El proyecto habilita referencias anulables con \texttt{<Nullable>enable</Nullable>}. Un valor que puede ser nulo se declara con \texttt{?}. En los puntos de entrada que reciben una referencia obligatoria desde una fuente externa o no validada, el parámetro se comprueba al inicio con \texttt{ArgumentNullException.ThrowIfNull}; la validación no se repite en métodos internos cuando el contrato y el análisis estático ya garantizan que el valor no es nulo. Las comparaciones se expresan con \texttt{is null} o \texttt{is not null}, y se prohíbe ocultar advertencias con el operador \texttt{!} sin una garantía demostrable (Microsoft, 2026f; Microsoft, 2026a).

\textbf{Código correcto:}
\begin{lstlisting}[style=csharpstyle]
public Player? FindPlayer(Guid playerId)
{
    return _players.SingleOrDefault(player => player.Id == playerId);
}

public void StartTurn(Player player)
{
    ArgumentNullException.ThrowIfNull(player);
    _turnManager.Start(player);
}
\end{lstlisting}

\textbf{Código incorrecto:}
\begin{lstlisting}[style=csharpstyle]
public Player FindPlayer(Guid playerId)
{
    return _players.SingleOrDefault(player => player.Id == playerId)!;
}

public void StartTurn(Player player)
{
    if (player == null)
    {
        throw new ArgumentNullException(nameof(player));
    }

    _turnManager.Start(player);
}
\end{lstlisting}
